\documentclass[14pt]{extarticle}
% =============================================================================
% Shared Preamble for Worksheets, Homeworks, Exams & Lesson Notes
% (scripts/tex_preamble.tex)
% =============================================================================

\newif\ifdetailed
\ifdefined\setdetailed
  \setdetailed
\fi

\newif\ifwithsols
\ifdefined\setwithsols
  \setwithsols
\fi

\newif\ifwhiteblack
\ifdefined\setwhiteblack
  \setwhiteblack
\fi

% Document-class specific packages (loaded before common_preamble.tex / polyglossia)
\usepackage{fancyhdr}
\usepackage{enumitem}
\usepackage{titlesec}
\usepackage[bottom]{footmisc}
\usepackage{titling}
\usepackage{longtable} % page-breakable tables
\usepackage{float}
\usepackage[onehalfspacing]{setspace}

\ifwhiteblack
  \usepackage[a4paper, landscape, top=1.25cm, bottom=1.25cm, left=1.25cm, right=1.25cm, headheight=1.5cm, headsep=0.5cm]{geometry}
\else
  \usepackage[a4paper, top=2.4cm, bottom=2.4cm, left=2.2cm, right=2.2cm, headheight=1.4cm, headsep=0.4cm, footskip=1.0cm]{geometry}
\fi

% Shared common packages, fonts, polyglossia, callout boxes, code helpers, date & header macros
\IfFileExists{scripts/common_preamble.tex}{%
  % =============================================================================
% Shared Common Preamble (scripts/common_preamble.tex)
% Included by:
%   - scripts/tex_preamble.tex    (worksheets, homeworks, exams, lesson notes)
%   - scripts/book_preamble.tex   (two-sided Arabic book)
%   - scripts/beamer_preamble.tex (Beamer presentation slides)
% =============================================================================

\usepackage{fontspec}
\usepackage{amsmath}
\usepackage{amssymb}
\usepackage{xcolor}
\usepackage{graphicx}
\usepackage{listings}
\usepackage[most]{tcolorbox}
\usepackage{pifont}
\usepackage{minted}
\usepackage{etoolbox}
\usepackage{array}
\usepackage{booktabs}
\usepackage{tikz}
\usepackage{datetime2}

\usetikzlibrary{automata, positioning, arrows.meta, arrows, shapes.geometric, calc}

% Beamer already loads hyperref internally; load it here for article/book classes
% right before polyglossia (bidi).
\makeatletter
\@ifclassloaded{beamer}{}{%
  \usepackage{hyperref}%
  \hypersetup{hidelinks}%
  \providecommand{\LTcaptype}{table}%
}
\makeatother

% Load polyglossia (and bidi) AFTER graphics, tables, boxes, code, and hyperref
\usepackage{polyglossia}
\setmainlanguage[numerals=western]{arabic}
\setotherlanguage{english}
\newfontfamily\arabicfont[Script=Arabic]{Calibri}
\newfontfamily\arabicfontsf[Script=Arabic]{Calibri}
\newfontfamily\arabicfonttt[Script=Arabic]{Courier New}
\newfontfamily\emoji[BoldFont=*,ItalicFont=*]{Segoe UI Emoji}

% Allow slight inter-word stretch to avoid Overfull \hbox with inline \textenglish
\emergencystretch=1.5em

% -----------------------------------------------------------------------------
% Default Minted & Callout Box Definitions
% -----------------------------------------------------------------------------
\AtBeginEnvironment{minted}{\setstretch{0.95}}
\setminted{fontsize=\small, breaklines=true, vspace=2pt}

\tcbset{
  arc=2.5pt, outer arc=2.5pt,
  boxsep=1pt, left=6pt, right=6pt, top=3pt, bottom=3pt,
  toptitle=1.5pt, bottomtitle=1.5pt,
  before skip=6pt, after skip=6pt,
  fonttitle=\bfseries,
  fontupper=\setstretch{1.1}\setlength{\parskip}{2pt},
  colback=white, colframe=black, boxrule=0.6pt,
  boxcolor/.style 2 args={colback=#1, colframe=#2}
}

\newtcolorbox{boxDef}[1][]{boxcolor={blue!4!white}{blue!75!black},
  title={{\emoji📘} تعريف\ifx\\#1\\\else ~#1\fi :}}
\newtcolorbox{boxExercise}[1][]{boxcolor={cyan!4!white}{cyan!70!black},
  title={{\emoji🧩} تمرين\ifx\\#1\\\else ~#1\fi :}}
\newtcolorbox{boxExample}[1][]{boxcolor={yellow!6!white}{orange!85!black},
  title={{\emoji📝} مثال\ifx\\#1\\\else ~#1\fi :}}
\newtcolorbox{boxNote}[1][]{boxcolor={gray!7!white}{black!75},
  title={{\emoji✨} ملاحظة\ifx\\#1\\\else ~#1\fi :}}
\newtcolorbox{boxAttention}[1][]{boxcolor={magenta!5!white}{magenta!75!black},
  title={{\emoji🔔} تنبيه\ifx\\#1\\\else ~#1\fi :}}
\newtcolorbox{boxWarning}[1][]{boxcolor={red!4!white}{red!75!black},
  title={{\emoji⚡} ملاحظة هامة\ifx\\#1\\\else ~#1\fi :}}
\newtcolorbox{boxSolution}[1][]{boxcolor={green!4!white}{green!60!black},
  title={{\emoji✅} حل\ifx\\#1\\\else ~#1\fi :}}
\newtcolorbox{boxSymbol}[1][]{boxcolor={purple!4!white}{purple!70!black},
  title={{\emoji🔣} رمز\ifx\\#1\\\else ~#1\fi :}}
\newtcolorbox{boxHint}[1][]{boxcolor={teal!4!white}{teal!65!black},
  title={{\emoji💡} تلميح\ifx\\#1\\\else ~#1\fi :}}

% Compact code boxes: tighter line spacing and small vertical gaps
\tcbset{
  simplecode/.style={
    arc=2pt, outer arc=2pt,
    left=4pt, right=4pt, top=4pt, bottom=4pt,
    fontupper=\setstretch{0.95},
    before skip=6pt, after skip=6pt,
    colback=gray!5, colframe=black!50, boxrule=0.4pt
  }
}
\newenvironment{boxCode}{\begin{tcolorbox}[simplecode]}{\end{tcolorbox}}

% Code cell used inside tables (keeps code left-to-right inside RTL tables)
% Use inside a fixed-width p{...} column: \begin{pycell}\begin{minted}...\end{minted}\end{pycell}
\makeatletter
\newtcolorbox{pycell}{frame empty,colback=white,boxsep=0pt,left=2pt,right=2pt,top=2pt,bottom=2pt,width=\linewidth,before=\noindent,after=,baseline={\tcb@height-\ht\strutbox},fontupper=\setstretch{0.95}\setlength{\topsep}{0pt}\setlength{\partopsep}{0pt}\setlength{\parskip}{0pt}}
\makeatother

\newcolumntype{C}[1]{>{\centering\arraybackslash}p{#1}}

% Dotted date style (DD.MM.YYYY)
\makeatletter
\DTMnewdatestyle{dotted}{%
  \renewcommand{\DTMdisplaydate}[4]{\two@digits{##3}.\two@digits{##2}.##1}%
  \renewcommand{\DTMDisplaydate}{\DTMdisplaydate}%
}
\makeatother
\DTMsetdatestyle{dotted}

% -----------------------------------------------------------------------------
% Customizable Header / Footer Metadata
% -----------------------------------------------------------------------------
\newcommand{\headerleft}{مدرسة التسامح الشاملة}
\newcommand{\headerright}{الأستاذ محمود اغبارية}
\newcommand{\footercenter}{\thepage}

\newcommand{\setheaderleft}[1]{\renewcommand{\headerleft}{#1}}
\newcommand{\setheaderright}[1]{\renewcommand{\headerright}{#1}}
\newcommand{\setfootercenter}[1]{\renewcommand{\footercenter}{#1}}
%
}{%
  \IfFileExists{../../../scripts/common_preamble.tex}{%
    % =============================================================================
% Shared Common Preamble (scripts/common_preamble.tex)
% Included by:
%   - scripts/tex_preamble.tex    (worksheets, homeworks, exams, lesson notes)
%   - scripts/book_preamble.tex   (two-sided Arabic book)
%   - scripts/beamer_preamble.tex (Beamer presentation slides)
% =============================================================================

\usepackage{fontspec}
\usepackage{amsmath}
\usepackage{amssymb}
\usepackage{xcolor}
\usepackage{graphicx}
\usepackage{listings}
\usepackage[most]{tcolorbox}
\usepackage{pifont}
\usepackage{minted}
\usepackage{etoolbox}
\usepackage{array}
\usepackage{booktabs}
\usepackage{tikz}
\usepackage{datetime2}

\usetikzlibrary{automata, positioning, arrows.meta, arrows, shapes.geometric, calc}

% Beamer already loads hyperref internally; load it here for article/book classes
% right before polyglossia (bidi).
\makeatletter
\@ifclassloaded{beamer}{}{%
  \usepackage{hyperref}%
  \hypersetup{hidelinks}%
  \providecommand{\LTcaptype}{table}%
}
\makeatother

% Load polyglossia (and bidi) AFTER graphics, tables, boxes, code, and hyperref
\usepackage{polyglossia}
\setmainlanguage[numerals=western]{arabic}
\setotherlanguage{english}
\newfontfamily\arabicfont[Script=Arabic]{Calibri}
\newfontfamily\arabicfontsf[Script=Arabic]{Calibri}
\newfontfamily\arabicfonttt[Script=Arabic]{Courier New}
\newfontfamily\emoji[BoldFont=*,ItalicFont=*]{Segoe UI Emoji}

% Allow slight inter-word stretch to avoid Overfull \hbox with inline \textenglish
\emergencystretch=1.5em

% -----------------------------------------------------------------------------
% Default Minted & Callout Box Definitions
% -----------------------------------------------------------------------------
\AtBeginEnvironment{minted}{\setstretch{0.95}}
\setminted{fontsize=\small, breaklines=true, vspace=2pt}

\tcbset{
  arc=2.5pt, outer arc=2.5pt,
  boxsep=1pt, left=6pt, right=6pt, top=3pt, bottom=3pt,
  toptitle=1.5pt, bottomtitle=1.5pt,
  before skip=6pt, after skip=6pt,
  fonttitle=\bfseries,
  fontupper=\setstretch{1.1}\setlength{\parskip}{2pt},
  colback=white, colframe=black, boxrule=0.6pt,
  boxcolor/.style 2 args={colback=#1, colframe=#2}
}

\newtcolorbox{boxDef}[1][]{boxcolor={blue!4!white}{blue!75!black},
  title={{\emoji📘} تعريف\ifx\\#1\\\else ~#1\fi :}}
\newtcolorbox{boxExercise}[1][]{boxcolor={cyan!4!white}{cyan!70!black},
  title={{\emoji🧩} تمرين\ifx\\#1\\\else ~#1\fi :}}
\newtcolorbox{boxExample}[1][]{boxcolor={yellow!6!white}{orange!85!black},
  title={{\emoji📝} مثال\ifx\\#1\\\else ~#1\fi :}}
\newtcolorbox{boxNote}[1][]{boxcolor={gray!7!white}{black!75},
  title={{\emoji✨} ملاحظة\ifx\\#1\\\else ~#1\fi :}}
\newtcolorbox{boxAttention}[1][]{boxcolor={magenta!5!white}{magenta!75!black},
  title={{\emoji🔔} تنبيه\ifx\\#1\\\else ~#1\fi :}}
\newtcolorbox{boxWarning}[1][]{boxcolor={red!4!white}{red!75!black},
  title={{\emoji⚡} ملاحظة هامة\ifx\\#1\\\else ~#1\fi :}}
\newtcolorbox{boxSolution}[1][]{boxcolor={green!4!white}{green!60!black},
  title={{\emoji✅} حل\ifx\\#1\\\else ~#1\fi :}}
\newtcolorbox{boxSymbol}[1][]{boxcolor={purple!4!white}{purple!70!black},
  title={{\emoji🔣} رمز\ifx\\#1\\\else ~#1\fi :}}
\newtcolorbox{boxHint}[1][]{boxcolor={teal!4!white}{teal!65!black},
  title={{\emoji💡} تلميح\ifx\\#1\\\else ~#1\fi :}}

% Compact code boxes: tighter line spacing and small vertical gaps
\tcbset{
  simplecode/.style={
    arc=2pt, outer arc=2pt,
    left=4pt, right=4pt, top=4pt, bottom=4pt,
    fontupper=\setstretch{0.95},
    before skip=6pt, after skip=6pt,
    colback=gray!5, colframe=black!50, boxrule=0.4pt
  }
}
\newenvironment{boxCode}{\begin{tcolorbox}[simplecode]}{\end{tcolorbox}}

% Code cell used inside tables (keeps code left-to-right inside RTL tables)
% Use inside a fixed-width p{...} column: \begin{pycell}\begin{minted}...\end{minted}\end{pycell}
\makeatletter
\newtcolorbox{pycell}{frame empty,colback=white,boxsep=0pt,left=2pt,right=2pt,top=2pt,bottom=2pt,width=\linewidth,before=\noindent,after=,baseline={\tcb@height-\ht\strutbox},fontupper=\setstretch{0.95}\setlength{\topsep}{0pt}\setlength{\partopsep}{0pt}\setlength{\parskip}{0pt}}
\makeatother

\newcolumntype{C}[1]{>{\centering\arraybackslash}p{#1}}

% Dotted date style (DD.MM.YYYY)
\makeatletter
\DTMnewdatestyle{dotted}{%
  \renewcommand{\DTMdisplaydate}[4]{\two@digits{##3}.\two@digits{##2}.##1}%
  \renewcommand{\DTMDisplaydate}{\DTMdisplaydate}%
}
\makeatother
\DTMsetdatestyle{dotted}

% -----------------------------------------------------------------------------
% Customizable Header / Footer Metadata
% -----------------------------------------------------------------------------
\newcommand{\headerleft}{مدرسة التسامح الشاملة}
\newcommand{\headerright}{الأستاذ محمود اغبارية}
\newcommand{\footercenter}{\thepage}

\newcommand{\setheaderleft}[1]{\renewcommand{\headerleft}{#1}}
\newcommand{\setheaderright}[1]{\renewcommand{\headerright}{#1}}
\newcommand{\setfootercenter}[1]{\renewcommand{\footercenter}{#1}}
%
  }{%
    \input{../../scripts/common_preamble.tex}%
  }%
}

\pagestyle{fancy}

% Prevent widow/orphan lines
\widowpenalty=10000
\clubpenalty=10000

% Compact section headings and default list spacing so pages do not waste vertical space
\titlespacing*{\section}{0pt}{1.4ex plus 0.4ex minus 0.2ex}{0.7ex plus 0.2ex}
\titlespacing*{\subsection}{0pt}{1.1ex plus 0.3ex minus 0.2ex}{0.5ex plus 0.1ex}
\titlespacing*{\subsubsection}{0pt}{0.9ex plus 0.2ex minus 0.1ex}{0.4ex plus 0.1ex}
\setlist{topsep=4pt, partopsep=0pt, parsep=2pt, itemsep=4pt}

\gappto\captionsarabic{\renewcommand{\contentsname}{جدول المحتويات}}

% Legacy listings settings
\lstset{
  language=[Sharp]C,
  numbers=left,
  stepnumber=1,
  numbersep=8pt,
  frame=single,
  basicstyle=\ttfamily\small
}
\ifwhiteblack
  \lstset{
    keywordstyle=\bfseries,
    stringstyle=\ttfamily,
    commentstyle=\itshape
  }
\else
  \lstset{
    keywordstyle=\color{blue},
    stringstyle=\color{red},
    commentstyle=\color{green!50!black}
  }
\fi

% White-and-black landscape presentation mode overrides
\ifwhiteblack
  \AtBeginDocument{
    \LARGE
    \usemintedstyle{bw}
  }
  \setminted{vspace=2pt}
  \tcbset{
    colback=white, colframe=black, boxrule=1pt,
    boxcolor/.style 2 args={colback=white, colframe=black},
    simplecode/.append style={colback=white, colframe=black, boxrule=1pt}
  }
\fi

% Default to empty author to avoid "LaTeX Warning: No \author given" when omitted
\author{}

% Compact title block
\makeatletter
\renewcommand{\@maketitle}{%
  \begin{center}%
    {\ifwhiteblack\Huge\else\LARGE\fi \bfseries \@title \par}%
    \ifdefempty{\@author}{}{%
      \vskip 0.15em%
      {\large \@author \par}%
    }%
    {\small \hfill \textenglish{\DTMtoday} \par}%
  \end{center}%
  \vspace{-0.3cm}%
}
\makeatother

\apptocmd{\maketitle}{\thispagestyle{fancy}}{}{}

% Apply page headers and footer (using \headerleft, \headerright, \footercenter from common_preamble.tex)
\newcommand{\applyfancyheaders}{%
  \fancyhf{}%
  \fancyhead[L]{\textarabic{\headerleft}}%
  \fancyhead[R]{\textarabic{\headerright}}%
  \fancyfoot[C]{\textarabic{\footercenter}}%
}

\fancypagestyle{plain}{\applyfancyheaders}
\applyfancyheaders

\setcounter{tocdepth}{3} % only section, subsection, and subsubsection in TOC

\newcommand{\importpdfpage}[2]{%
    \newgeometry{left=0cm,right=0cm,top=0cm,bottom=0cm}
    \thispagestyle{empty}
    \includegraphics[page=#2,width=0.95\paperwidth,height=0.95\paperheight,keepaspectratio]{#1}
    \restoregeometry
}

\newcommand{\insertFullImg}[1]{%
    \noindent
    \makebox[\textwidth][c]{
        \includegraphics[width=0.9\paperwidth,keepaspectratio]{#1}%
    }%
}%

\setheaderright{بايتون - السلاسل النصية}

\ifwithsols
\title{حل ورقة عمل - السلاسل النصية \textenglish{Strings} (بايثون أ)}
\else
\title{ورقة عمل - السلاسل النصية \textenglish{Strings} (بايثون أ)}
\fi

\begin{document}

\maketitle
\thispagestyle{fancy}

\section*{القسم الأول: أسئلة الدرس - السلاسل النصية}

\begin{enumerate}[label=\textbf{تمرين \arabic*:}, wide=0pt, itemsep=2.2em]

% ---------------------------------------------------------------
\item أمامكم مقطع البرنامج التالي الذي يستعمل العملية \textenglish{count}:

\begin{boxCode}
\begin{english}
\begin{minted}{python}
text = "I love bananas, banana are my favorite fruit"
x = text.count("banana")
print("x=", x)
\end{minted}
\end{english}
\end{boxCode}

\begin{enumerate}[label=\alph*.]
\item ما هو المُخرج الناتج من تشغيل مقطع البرنامج؟
\ifwithsols
\begin{boxSolution}
المُخرج الناتج هو: \textenglish{\texttt{x= 2}} (لأنّ السلسلة النصية الجزئية \textenglish{\texttt{"banana"}} تظهر مرتين: مرّة داخل الكلمة \textenglish{\texttt{"bananas"}} ومرّة ككلمة مستقلة \textenglish{\texttt{"banana"}}).
\end{boxSolution}
\fi
\filbreak

\item انسخوا الكود الذي أمامكم وبدّلوا السطر الثالث بالسطر التالي:
\begin{boxCode}
\begin{english}
\begin{minted}{python}
print(f"count banana in text is : {x}")
\end{minted}
\end{english}
\end{boxCode}
ما هو المُخرج الناتج الآن؟
\ifwithsols
\begin{boxSolution}
المُخرج الناتج هو: \textenglish{\texttt{count banana in text is : 2}}
\end{boxSolution}
\fi
\filbreak
\end{enumerate}
\filbreak

% ---------------------------------------------------------------
\item معطاة السلسلة النصية التالية (طولها $33$ رمزًا، مرقّمة بالمؤشرات الموجبة من $0$ إلى $32$، وبالمؤشرات السالبة من $-33$ إلى $-1$):

\begin{boxCode}
\begin{english}
\begin{minted}{python}
text = "think big be bold show your color"
\end{minted}
\end{english}
\end{boxCode}

ما هي قيمة الحدّ \textenglish{\texttt{text[16]}}؟ وما هي قيمة الحدّ \textenglish{\texttt{text[-17]}}؟

\ifwithsols
\begin{boxSolution}
قيمة الحدّ \textenglish{\texttt{text[16]}} وقيمة الحدّ \textenglish{\texttt{text[-17]}} هي الرمز \textenglish{\texttt{'d'}} (الحرف الأخير في الكلمة \textenglish{\texttt{"bold"}})، لأنّ المؤشر $16$ من البداية والمؤشر $-17$ من النهاية يشيران إلى نفس الموقع في السلسلة النصية ($16 + 17 = 33$ وهو طول السلسلة النصية).
\end{boxSolution}
\fi
\filbreak

% ---------------------------------------------------------------
\item معطاة السلسلة النصية التالية:

\begin{boxCode}
\begin{english}
\begin{minted}{python}
text = "sometimes you win sometimes you learn"
\end{minted}
\end{english}
\end{boxCode}

\begin{enumerate}[label=\alph*.]
\item استعينوا بالعملية \textenglish{find} لكي تعيدوا موقع ظهور السلسلة النصية الجزئية \textenglish{\texttt{"some"}} في السلسلة النصية \textenglish{text} وتطبعوه. ما هو المُخرج الناتج؟
\ifwithsols
\begin{boxSolution}
\begin{english}
\begin{minted}{python}
x = text.find("some")
print(x)
\end{minted}
\end{english}
المُخرج الناتج هو: $0$.
\end{boxSolution}
\fi
\filbreak

\item اشرحوا لماذا قيمة المتغير $x$ تساوي $0$ رغم أنّ السلسلة النصية الجزئية \textenglish{\texttt{"some"}} تظهر أكثر من مرّة في السلسلة النصية \textenglish{text}.
\ifwithsols
\begin{boxSolution}
لأنّ العملية \textenglish{find} تعيد مؤشر (موقع) \textbf{الظهور الأول} للسلسلة النصية الجزئية من جهة اليسار، وبما أنّ السلسلة النصية \textenglish{text} تبدأ بالكلمة \textenglish{\texttt{"sometimes"}} فإنّ أول ظهور لـ \textenglish{\texttt{"some"}} يبدأ في المؤشر $0$.
\end{boxSolution}
\fi
\filbreak
\end{enumerate}
\filbreak

\end{enumerate}

\clearpage
\section*{القسم الثاني: ورقة عمل - نتدرّب على حل المشكلات}

\begin{enumerate}[resume, label=\textbf{تمرين \arabic*:}, wide=0pt, itemsep=2.2em]

% ---------------------------------------------------------------
% سؤال 1 في ورقة العمل
\item \textbf{عمليات وصل ($+$) ومضاعفة ($*$):}
\begin{enumerate}[label=\alph*.]
\item اكتبوا المُخرج الناتج من مقطع البرنامج التالي:

\begin{boxCode}
\begin{english}
\begin{minted}{python}
text1 = "I like to play, "
text2 = "piano"
text3 = "guitar"
print(text1 + text2)
print(text1 + text3)
print(text3 * 2)
print(text2 * 3)
\end{minted}
\end{english}
\end{boxCode}

\ifwithsols
\begin{boxSolution}
\begin{english}
\begin{minted}{text}
I like to play, piano
I like to play, guitar
guitarguitar
pianopianopiano
\end{minted}
\end{english}
\end{boxSolution}
\fi
\filbreak

\item ما هي وظيفة العامل $+$ بين سلسلتين نصّيتين؟
\ifwithsols
\begin{boxSolution}
يقوم بدمج (وصل) السلسلتين النصّيتين معًا لتكوين سلسلة نصية واحدة جديدة.
\end{boxSolution}
\fi
\filbreak

\item ما هي وظيفة العامل $*$ بين سلسلة نصية وعدد صحيح موجب؟
\ifwithsols
\begin{boxSolution}
يقوم بتكرار (مضاعفة) السلسلة النصية بعدد المرات المحدّد بالعدد الصحيح.
\end{boxSolution}
\fi
\filbreak
\end{enumerate}
\filbreak

% ---------------------------------------------------------------
% سؤال 2 في ورقة العمل
\item
\begin{enumerate}[label=\alph*.]
\item اكتبوا المُخرج الناتج من مقطع البرنامج التالي:

\begin{boxCode}
\begin{english}
\begin{minted}{python}
num1 = "123"
num2 = "45"
num3 = "7"
print(num1 + num2)
print(num2 + num3)
print(num2 * 4)
print(num3 * 9)
\end{minted}
\end{english}
\end{boxCode}

\ifwithsols
\begin{boxSolution}
\begin{english}
\begin{minted}{text}
12345
457
45454545
777777777
\end{minted}
\end{english}
\end{boxSolution}
\fi
\filbreak

\item استعملوا العاملين $+$ و $*$ والمتغيرين \textenglish{num2} و \textenglish{num3} فقط، واكتبوا أمر طباعة ينتج عنه المُخرج \textenglish{\texttt{457457457}}.
\ifwithsols
\begin{boxSolution}
\begin{english}
\begin{minted}{python}
print((num2 + num3) * 3)
\end{minted}
\end{english}
\end{boxSolution}
\fi
\filbreak
\end{enumerate}
\filbreak

% ---------------------------------------------------------------
% سؤال 3 في ورقة العمل
\item \textbf{طول السلسلة النصية | \textenglish{len}:}
\begin{enumerate}[label=\alph*.]
\item اكتبوا المُخرج الناتج من مقطع البرنامج التالي:

\begin{boxCode}
\begin{english}
\begin{minted}{python}
text1 = "I like to play, "
text2 = "piano"
text3 = "guitar"
print(len(text1))
print(len(text2))
print(len(text2 + text3))
print(len(text3 + text2))
\end{minted}
\end{english}
\end{boxCode}

\ifwithsols
\begin{boxSolution}
\begin{english}
\begin{minted}{text}
16
5
11
11
\end{minted}
\end{english}
\end{boxSolution}
\fi
\filbreak

\item بماذا تقوم العملية \textenglish{len}؟
\ifwithsols
\begin{boxSolution}
تعيد طول السلسلة النصية، أي عدد جميع الرموز فيها (بما يشمل الحروف، الأرقام، الفراغات، وعلامات الترقيم).
\end{boxSolution}
\fi
\filbreak
\end{enumerate}
\filbreak

% ---------------------------------------------------------------
% سؤال 4 في ورقة العمل
\item \textbf{تحويل بين الأنماط | \textenglish{str} و \textenglish{int}:}
أمامكم مقطع برنامج فيه متغيران: \textenglish{num\_1} يحصل على السلسلة النصية \textenglish{\texttt{"123"}}، و \textenglish{num\_2} يحصل على العدد الصحيح $451$:

\begin{boxCode}
\begin{english}
\begin{minted}{python}
num_1 = "123"
num_2 = 451
print(num_1)
print(int(num_1))
print(num_1 + str(num_2))
print(int(num_1) + num_2)
word = num_1 + str(num_2)
print(word[-4])
\end{minted}
\end{english}
\end{boxCode}

\begin{enumerate}[label=\alph*.]
\item اكتبوا المُخرج الناتج من مقطع البرنامج أعلاه.
\ifwithsols
\begin{boxSolution}
\begin{english}
\begin{minted}{text}
123
123
123451
574
3
\end{minted}
\end{english}
(توضيح السطر الأخير: قيمة \textenglish{word} هي \textenglish{\texttt{"123451"}}، والرمز في المؤشر $-4$ من النهاية هو \textenglish{\texttt{'3'}}.)
\end{boxSolution}
\fi
\filbreak

\item بماذا تقوم العملية \textenglish{int()}؟
\ifwithsols
\begin{boxSolution}
تحوّل القيمة المعطاة (مثل سلسلة نصية تمثّل عددًا صحيحًا أو عددًا عشريًا) إلى عدد صحيح من نمط \textenglish{int}.
\end{boxSolution}
\fi
\filbreak

\item بماذا تقوم العملية \textenglish{str()}؟
\ifwithsols
\begin{boxSolution}
تحوّل القيمة المعطاة (مثل عدد صحيح أو عشري) إلى سلسلة نصية من نمط \textenglish{str}.
\end{boxSolution}
\fi
\filbreak
\end{enumerate}
\filbreak

% ---------------------------------------------------------------
% سؤال 9 في ورقة العمل
\item \textbf{الوصول إلى حدّ بواسطة المؤشر \textenglish{index[]}:}
معطى المتغير \textenglish{\texttt{word = "piano"}}. اكتبوا بجانب كل أمر طباعة المُخرج الناتج منه:

{\renewcommand{\arraystretch}{1.4}
\begin{longtable}{|p{6.5cm}|C{5cm}|}
\hline
\centering\textbf{الأمر بلغة بايثون} & \textbf{المُخرج الناتج} \\
\hline
\textenglish{\texttt{print(word[0])}} & \ifwithsols \textenglish{\texttt{p}} \fi \\ \hline
\textenglish{\texttt{print(word[1])}} & \ifwithsols \textenglish{\texttt{i}} \fi \\ \hline
\textenglish{\texttt{print(word[4])}} & \ifwithsols \textenglish{\texttt{o}} \fi \\ \hline
\textenglish{\texttt{print(word[-1])}} & \ifwithsols \textenglish{\texttt{o}} \fi \\ \hline
\textenglish{\texttt{print(word[-2])}} & \ifwithsols \textenglish{\texttt{n}} \fi \\ \hline
\textenglish{\texttt{print(word[-3])}} & \ifwithsols \textenglish{\texttt{a}} \fi \\ \hline
\textenglish{\texttt{print(word[-4])}} & \ifwithsols \textenglish{\texttt{i}} \fi \\ \hline
\end{longtable}}
\filbreak

% ---------------------------------------------------------------
% سؤال 10 في ورقة العمل
\item بالاعتماد على المتغير \textenglish{\texttt{word = "piano"}} من التمرين السابق، ما هو المُخرج الناتج من تنفيذ الأمر \textenglish{\texttt{print(word[5])}}؟ ولماذا؟

\ifwithsols
\begin{boxSolution}
ينتج خطأ وقت تشغيل (\textenglish{IndexError: string index out of range}) لأنّ طول السلسلة النصية \textenglish{\texttt{"piano"}} هو $5$، فالمؤشرات الموجبة تبدأ من $0$ وتنتهي عند $4$، ولا يوجد رمز في المؤشر $5$.
\end{boxSolution}
\fi
\filbreak

% ---------------------------------------------------------------
% سؤال 11 في ورقة العمل
\item معطى المتغير \textenglish{\texttt{st = "ABCDE"}}. أمامكم جدول له عمودان: لكل قيمة مُخرج مطلوبة في الجدول، اكتبوا أمرين مختلفين بلغة بايثون (أحدهما بمؤشر موجب والآخر بمؤشر سالب) بحيث يعطي كلٌّ منهما المُخرج المعطى (استعينوا بالبند الأول كمثال):

{\renewcommand{\arraystretch}{1.4}
\begin{longtable}{|p{8cm}|C{3.5cm}|}
\hline
\centering\textbf{الأمران في لغة بايثون} & \textbf{المُخرج} \\
\hline
\begin{pycell}
\begin{english}
\begin{minted}{python}
print(st[4])
print(st[-1])
\end{minted}
\end{english}
\end{pycell}
& \textenglish{\texttt{E}} \\
\hline
\ifwithsols
\begin{pycell}
\begin{english}
\begin{minted}{python}
print(st[3])
print(st[-2])
\end{minted}
\end{english}
\end{pycell}
\else \rule{0pt}{1.4cm} \fi
& \textenglish{\texttt{D}} \\
\hline
\ifwithsols
\begin{pycell}
\begin{english}
\begin{minted}{python}
print(st[2])
print(st[-3])
\end{minted}
\end{english}
\end{pycell}
\else \rule{0pt}{1.4cm} \fi
& \textenglish{\texttt{C}} \\
\hline
\ifwithsols
\begin{pycell}
\begin{english}
\begin{minted}{python}
print(st[1])
print(st[-4])
\end{minted}
\end{english}
\end{pycell}
\else \rule{0pt}{1.4cm} \fi
& \textenglish{\texttt{B}} \\
\hline
\ifwithsols
\begin{pycell}
\begin{english}
\begin{minted}{python}
print(st[0])
print(st[-5])
\end{minted}
\end{english}
\end{pycell}
\else \rule{0pt}{1.4cm} \fi
& \textenglish{\texttt{A}} \\
\hline
\end{longtable}}
\filbreak

% ---------------------------------------------------------------
% سؤال 12 في ورقة العمل
\item \textbf{اقتطاع السلاسل النصية \textenglish{text[start:end:jump]}:}
معطى المتغير التالي:
\begin{boxCode}
\begin{english}
\begin{minted}{python}
text = "I like to play drums"
\end{minted}
\end{english}
\end{boxCode}
أكملوا الجدول التالي بكتابة المُخرج الناتج لكل أمر، أو بكتابة أمر طباعة ملائم يعطي المُخرج المسجّل:

{\renewcommand{\arraystretch}{1.4}
\begin{longtable}{|p{7.5cm}|C{6cm}|}
\hline
\centering\textbf{الأمر في لغة بايثون} & \textbf{المُخرج} \\
\hline
\textenglish{\texttt{print(text[2:6])}} & \ifwithsols \textenglish{\texttt{like}} \fi \\ \hline
\textenglish{\texttt{print(text[6:])}} & \ifwithsols \textenglish{\texttt{\ to play drums}} \fi \\ \hline
\textenglish{\texttt{print(text[:15])}} & \ifwithsols \textenglish{\texttt{I like to play\ }} \fi \\ \hline
\ifwithsols \textenglish{\texttt{print(text[15:])}} \fi & \textenglish{\texttt{drums}} \\ \hline
\ifwithsols \textenglish{\texttt{print(text[7:14])}} \fi & \textenglish{\texttt{to play}} \\ \hline
\textenglish{\texttt{print(text[2:20:2])}} & \ifwithsols \textenglish{\texttt{lk opa rm}} \fi \\ \hline
\textenglish{\texttt{print(text[1:15:3])}} & \ifwithsols \textenglish{\texttt{\ ktpy}} \fi \\ \hline
\textenglish{\texttt{print(text[::-1])}} & \ifwithsols \textenglish{\texttt{smurd yalp ot ekil I}} \fi \\ \hline
\ifwithsols \textenglish{\texttt{print(text[0:17:4])}} \fi & \textenglish{\texttt{Ikoar}} \\ \hline
\ifwithsols \textenglish{\texttt{print(text[4:9:4])}} \fi & \textenglish{\texttt{ko}} \\ \hline
\textenglish{\texttt{print(text[-10:-6])}} & \ifwithsols \textenglish{\texttt{play}} \fi \\ \hline
\textenglish{\texttt{print(text[-20:-14])}} & \ifwithsols \textenglish{\texttt{I like}} \fi \\ \hline
\textenglish{\texttt{print(text[-20:-14:2])}} & \ifwithsols \textenglish{\texttt{Ilk}} \fi \\ \hline
\end{longtable}}
\filbreak

% ---------------------------------------------------------------
% سؤال 13 في ورقة العمل
\item معطى المتغير التالي:
\begin{boxCode}
\begin{english}
\begin{minted}{python}
number = "572483291"
\end{minted}
\end{english}
\end{boxCode}
لكل قيمة مُخرج في الجدول، اكتبوا أمرين مختلفين (أحدهما بمؤشرات موجبة والآخر بمؤشرات سالبة) بحيث ينتج عن كلّ أمر منهما المُخرج المطلوب (استعينوا بالبند الأول كمثال):

{\renewcommand{\arraystretch}{1.4}
\begin{longtable}{|p{8.5cm}|C{4cm}|}
\hline
\centering\textbf{الأمران في لغة بايثون} & \textbf{المُخرج} \\
\hline
\begin{pycell}
\begin{english}
\begin{minted}{python}
print(number[0:3])
print(number[-9:-6])
\end{minted}
\end{english}
\end{pycell}
& \textenglish{\texttt{572}} \\
\hline
\ifwithsols
\begin{pycell}
\begin{english}
\begin{minted}{python}
print(number[3:6])
print(number[-6:-3])
\end{minted}
\end{english}
\end{pycell}
\else \rule{0pt}{1.4cm} \fi
& \textenglish{\texttt{483}} \\
\hline
\ifwithsols
\begin{pycell}
\begin{english}
\begin{minted}{python}
print(number[6:9])
print(number[-3:])
\end{minted}
\end{english}
\end{pycell}
\else \rule{0pt}{1.4cm} \fi
& \textenglish{\texttt{291}} \\
\hline
\ifwithsols
\begin{pycell}
\begin{english}
\begin{minted}{python}
print(number[0:9:2])
print(number[-9::2])
\end{minted}
\end{english}
\end{pycell}
\else \rule{0pt}{1.4cm} \fi
& \textenglish{\texttt{52821}} \\
\hline
\ifwithsols
\begin{pycell}
\begin{english}
\begin{minted}{python}
print(number[3:7:3])
print(number[-6:-2:3])
\end{minted}
\end{english}
\end{pycell}
\else \rule{0pt}{1.4cm} \fi
& \textenglish{\texttt{42}} \\
\hline
\ifwithsols
\begin{pycell}
\begin{english}
\begin{minted}{python}
print(number[1:8:3])
print(number[-8:-1:3])
\end{minted}
\end{english}
\end{pycell}
\else \rule{0pt}{1.4cm} \fi
& \textenglish{\texttt{789}} \\
\hline
\end{longtable}}
\filbreak

% ---------------------------------------------------------------
% سؤال 14 في ورقة العمل
\item \textbf{العامل المنطقي \textenglish{in}:}
\begin{enumerate}[label=\alph*.]
\item اكتبوا المُخرج الناتج من مقطع البرنامج التالي:

\begin{boxCode}
\begin{english}
\begin{minted}{python}
print('a' in 'jerusalem')
print('x' in 'jerusalem')
print('2' in str(12345))
print('9' in str(12345))
\end{minted}
\end{english}
\end{boxCode}

\ifwithsols
\begin{boxSolution}
\begin{english}
\begin{minted}{text}
True
False
True
False
\end{minted}
\end{english}
\end{boxSolution}
\fi
\filbreak

\item ماذا ينفّذ العامل \textenglish{in} بين سلسلتين نصّيتين؟
\ifwithsols
\begin{boxSolution}
يفحص ما إذا كانت السلسلة النصية الجزئية (على يسار \textenglish{in}) موجودة داخل السلسلة النصية الأصلية (على يمين \textenglish{in})، ويعيد \textenglish{True} إذا وُجدت أو \textenglish{False} إذا لم توجد.
\end{boxSolution}
\fi
\filbreak

\item اكتبوا مقطع برنامج يستقبل عددًا صحيحًا، ويفحص بواسطة العامل \textenglish{in} ما إذا كانت المنزلة $9$ أو المنزلة $3$ تظهر في العدد، ويطبع \textenglish{True} أو \textenglish{False}.
\ifwithsols
\begin{boxSolution}
\begin{english}
\begin{minted}{python}
num = int(input("Enter an integer: "))
s = str(num)
print('9' in s or '3' in s)
\end{minted}
\end{english}
\end{boxSolution}
\fi
\filbreak
\end{enumerate}
\filbreak

% ---------------------------------------------------------------
% سؤال 15 في ورقة العمل
\item \textbf{العملية \textenglish{find}:}
\begin{enumerate}[label=\alph*.]
\item اكتبوا المُخرج الناتج من مقطع البرنامج التالي:

\begin{boxCode}
\begin{english}
\begin{minted}{python}
text = "I LIKE TO PLAY DRUMS"
print(text.find('L'))
print(text.find('LIKE'))
print(text.find('PIANO'))
print(text.find('L', 3))
print(text.find('L', 12))
\end{minted}
\end{english}
\end{boxCode}

\ifwithsols
\begin{boxSolution}
\begin{english}
\begin{minted}{text}
2
2
-1
11
-1
\end{minted}
\end{english}
\end{boxSolution}
\fi
\filbreak

\item ماذا تنفّذ العملية \textenglish{find(sub, begin, end)}؟
\ifwithsols
\begin{boxSolution}
تبحث عن السلسلة النصية الجزئية \textenglish{sub} داخل السلسلة النصية الأصلية بدءًا من المؤشر \textenglish{begin} وحتى المؤشر \textenglish{end} (لا يشمل \textenglish{end})، وتعيد مؤشر أول ظهور لها، أو تعيد $-1$ إذا لم توجد في هذا المجال.
\end{boxSolution}
\fi
\filbreak
\end{enumerate}
\filbreak

% ---------------------------------------------------------------
% سؤال 16 في ورقة العمل
\item اكتبوا مقطع برنامج يستقبل سلسلة نصية من المستخدم. إذا كانت السلسلة النصية تحتوي على الكلمة \textenglish{\texttt{"war"}}، يجب تبديلها بالكلمة \textenglish{\texttt{"flower"}}، ثم طباعة السلسلة النصية الناتجة.

\ifwithsols
\begin{boxSolution}
\begin{english}
\begin{minted}{python}
text = input("Enter a string: ")
if "war" in text:
    text = text.replace("war", "flower")
print(text)
\end{minted}
\end{english}
\end{boxSolution}
\fi
\filbreak

% ---------------------------------------------------------------
% سؤال 17 في ورقة العمل
\item \textbf{العمليتان \textenglish{lower()} و \textenglish{upper()}:}
\begin{enumerate}[label=\alph*.]
\item اكتبوا المُخرج الناتج من مقطع البرنامج التالي:

\begin{boxCode}
\begin{english}
\begin{minted}{python}
print("piano".upper())
print("DRUMS".lower())
print("Piano And Drums".upper())
\end{minted}
\end{english}
\end{boxCode}

\ifwithsols
\begin{boxSolution}
\begin{english}
\begin{minted}{text}
PIANO
drums
PIANO AND DRUMS
\end{minted}
\end{english}
\end{boxSolution}
\fi
\filbreak

\item ماذا تنفّذ كل واحدة من العمليتين التاليتين؟
\begin{itemize}
\item العملية \textenglish{upper()}:
\item العملية \textenglish{lower()}:
\end{itemize}
\ifwithsols
\begin{boxSolution}
\begin{itemize}
\item \textenglish{upper()}: تعيد سلسلة نصية جديدة تمّ فيها تحويل جميع الأحرف الإنجليزية إلى أحرف كبيرة (\textenglish{Capital letters}) دون تغيير السلسلة النصية الأصلية.
\item \textenglish{lower()}: تعيد سلسلة نصية جديدة تمّ فيها تحويل جميع الأحرف الإنجليزية إلى أحرف صغيرة (\textenglish{Small letters}) دون تغيير السلسلة النصية الأصلية.
\end{itemize}
\end{boxSolution}
\fi
\filbreak
\end{enumerate}
\filbreak

% ---------------------------------------------------------------
% سؤال 18 في ورقة العمل
\item اكتبوا مقطع برنامج يستقبل سلسلة نصية من المستخدم، ويطبع \textenglish{True} إذا كانت السلسلة النصية تحتوي على الكلمة \textenglish{\texttt{"yes"}} أو \textenglish{\texttt{"YES"}}، وخلاف ذلك يطبع \textenglish{False}. استعملوا إحدى العمليتين \textenglish{upper()} أو \textenglish{lower()}.

\ifwithsols
\begin{boxSolution}
\begin{english}
\begin{minted}{python}
text = input("Enter a string: ")
print("yes" in text.lower())
\end{minted}
\end{english}
\end{boxSolution}
\fi
\filbreak

% ---------------------------------------------------------------
% سؤال 19 في ورقة العمل
\item \textbf{العمليتان \textenglish{startswith()} و \textenglish{endswith()}:}
\begin{enumerate}[label=\alph*.]
\item اكتبوا المُخرج الناتج من مقطع البرنامج التالي:

\begin{boxCode}
\begin{english}
\begin{minted}{python}
text = "I like to play drums"
print(text.startswith('I like'))
print(text.startswith('play', 10, 17))
print(text.startswith('to play', 10, 17))
\end{minted}
\end{english}
\end{boxCode}

\ifwithsols
\begin{boxSolution}
\begin{english}
\begin{minted}{text}
True
True
False
\end{minted}
\end{english}
\end{boxSolution}
\fi
\filbreak

\item اشرحوا باختصار مبنى وهدف كلٍّ من العمليتين التاليتين:
\begin{itemize}
\item \textenglish{startswith(sub, begin, end)}:
\item \textenglish{endswith(sub, begin, end)}:
\end{itemize}
\ifwithsols
\begin{boxSolution}
\begin{itemize}
\item \textenglish{startswith(sub, begin, end)}: تفحص ما إذا كان المقطع الجزئي من السلسلة النصية بين المؤشر \textenglish{begin} والمؤشر \textenglish{end} (لا يشمل \textenglish{end}) \textbf{يبدأ} بالسلسلة النصية الجزئية \textenglish{sub}، وتعيد \textenglish{True} أو \textenglish{False}.
\item \textenglish{endswith(sub, begin, end)}: تفحص ما إذا كان المقطع الجزئي من السلسلة النصية بين المؤشر \textenglish{begin} والمؤشر \textenglish{end} (لا يشمل \textenglish{end}) \textbf{ينتهي} بالسلسلة النصية الجزئية \textenglish{sub}، وتعيد \textenglish{True} أو \textenglish{False}.
\end{itemize}
\end{boxSolution}
\fi
\filbreak
\end{enumerate}
\filbreak

% ---------------------------------------------------------------
% سؤال 20 في ورقة العمل
\item \textbf{العملية \textenglish{replace}:}
اشرحوا مبنى وهدف العملية \textenglish{replace(old, new, max)}.

\ifwithsols
\begin{boxSolution}
تعيد سلسلة نصية جديدة بعد استبدال السلسلة النصية الجزئية \textenglish{old} بالسلسلة النصية الجزئية \textenglish{new} بعدد مرات أقصاه \textenglish{max} من بداية السلسلة النصية (دون تغيير السلسلة النصية الأصلية).
\end{boxSolution}
\fi
\filbreak

% ---------------------------------------------------------------
% سؤال 21 في ورقة العمل
\item اكتبوا مقطع برنامج يستقبل سلسلة نصية، ويستبدل الرمز \textenglish{\texttt{"|"}} بالسلسلة النصية \textenglish{\texttt{"we"}} في المرات الثلاث الأولى فقط التي يظهر فيها الرمز \textenglish{\texttt{"|"}}، ثم يطبع السلسلة النصية الناتجة.

\ifwithsols
\begin{boxSolution}
\begin{english}
\begin{minted}{python}
text = input("Enter a string: ")
new_text = text.replace("|", "we", 3)
print(new_text)
\end{minted}
\end{english}
\end{boxSolution}
\fi
\filbreak

% ---------------------------------------------------------------
% سؤال 22 في ورقة العمل
\item \textbf{العملية \textenglish{count}:}
اكتبوا المُخرج الناتج من مقطع البرنامج التالي:

\begin{boxCode}
\begin{english}
\begin{minted}{python}
text = "I like to play guitar and piano and drums"
print(text.count("and"))
print(text.count("and", 25))
print(text.count("play"))
print(text.count("books"))
\end{minted}
\end{english}
\end{boxCode}

\ifwithsols
\begin{boxSolution}
\begin{english}
\begin{minted}{text}
2
1
1
0
\end{minted}
\end{english}
\end{boxSolution}
\fi
\filbreak

% ---------------------------------------------------------------
% سؤال 23 في ورقة العمل
\item اشرحوا مبنى وهدف العملية \textenglish{count(sub, begin, end)}.

\ifwithsols
\begin{boxSolution}
تحصي وتعيد عدد المرات التي تظهر فيها السلسلة النصية الجزئية \textenglish{sub} داخل السلسلة النصية الأصلية في المجال الممتد من المؤشر \textenglish{begin} وحتى المؤشر \textenglish{end} (لا يشمل \textenglish{end}).
\end{boxSolution}
\fi
\filbreak

% ---------------------------------------------------------------
% سؤال 24 في ورقة العمل
\item اكتبوا مقطع برنامج يستقبل سلسلة نصية من المستخدم، ويحصي ويطبع عدد المرات التي تظهر فيها السلسلة النصية الجزئية \textenglish{\texttt{"the"}} في السلسلة النصية.

\ifwithsols
\begin{boxSolution}
\begin{english}
\begin{minted}{python}
text = input("Enter a string: ")
print(text.count("the"))
\end{minted}
\end{english}
\end{boxSolution}
\fi
\filbreak

\end{enumerate}

\vspace{1cm}
\begin{flushleft}
نرجو لكم وقتًا ممتعًا.

معلمو الموضوع.
\end{flushleft}

\end{document}
